\begin{helper}[Transport normal data along an exact-trunk candidate]
Let $N$ be normal globally maximum data for $(T,e_0,\varphi)$, and let
$S$ be a \noderef{ScheideBalancedCandidate} for $N$.  If
$\noderef{TreeSequenceVertexSet}(S)=
\noderef{TreeSequenceVertexSet}(T)$, then there are normal globally maximum
data $N'$ for $(S,e_0,\varphi)$ with
$N'.\operatorname{height}=N.\operatorname{height}$.
\end{helper}
\begin{proof}
Put $p=N.\operatorname{height}$.  We first establish equality of the
displayed lengths, which is needed to transport all bounded indices.
For any \noderef{TreeSequence} $R$, its vertex map is injective and
\noderef{TreeSequenceVertexSet} is exactly the image of that map on
$\operatorname{Fin}(R.\operatorname{length}+1)$.  Therefore
$|V(R)|=R.\operatorname{length}+1$.  The assumed equality of vertex sets
gives
\[
 S.\operatorname{length}+1=|V(S)|=|V(T)|
 =T.\operatorname{length}+1,
\]
and cancellation yields $S.\operatorname{length}=T.\operatorname{length}$.

Next unpack the exact-trunk field of
\noderef{ScheideBalancedCandidate}.  By
\noderef{NormalTripleTrunk}, the trunk is the prefix of $T$ of length
$p-1$; hence its vertices in positions $0,\ldots,p-1$ and edges in
positions $0,\ldots,p-2$ are literally those of $T$.  The
\noderef{TreeSequenceInitialSegment} field says, entry by entry and with
the corresponding finite-index casts, that this same trunk is the initial
segment of $S$.  Consequently
\[
 S.\operatorname{vertex}(a)=T.\operatorname{vertex}(a)
 \quad(0\le a<p),
 \qquad
 S.\operatorname{edge}(b)=T.\operatorname{edge}(b)
 \quad(0\le b<p-1).
\]
These are ordered equalities of exact vertex and edge identities, not merely
consequences of equality of the two vertex sets.

We now construct every field of $N'$ from the fields of
\noderef{NormalGlobalMaximumTriple}.  Its global field, in the sense of
\noderef{GlobalMaximumTriple}, uses the original $k$-triple, the
candidate's Tashkinov field, and
\[
 |V(S)|=|V(T)|=\noderef{TashkinovOrder}(G),
\]
where the last equality is the global field of $N$.  Set the new height to
$p$ and retain $N$'s lower bound.  Its upper bound transports along the
proved equality of lengths.  Retain the same colors $\alpha,\beta$.
Their missingness at the first two vertices transports along the exact
prefix equalities just proved.

Retain the same closing edge identity.  Transport $N$'s last index across
$S.\operatorname{length}=T.\operatorname{length}$; its value remains
$p-1$.  The transported vertex at that index and the vertex at index zero
agree with their counterparts in $T$, so the closing-edge endpoint equation
is unchanged.  Retain the same closing incidence path: its ambient
two-color incidence graph depends only on $G,\varphi,\alpha,\beta$, none of
which changed.  Its simplicity and length are therefore unchanged.

For the path-support fields, the vertex-node equivalence in $N$ says that
the supported vertices are exactly the first $p$ displayed vertices of
$T$.  The exact-prefix equalities replace each of these, at the same
bounded index, by the corresponding vertex of $S$, proving the new
vertex-support equivalence in both directions.  The edge-node equivalence
uses the same closing identity together with the exact trunk edges
$e_2,\ldots,e_{p-1}$, equivalently the zero-based positions $b$ satisfying
$1\le b$ and $b+2\le p$.  Each such edge of $S$ equals the identically
indexed edge of $T$, so both directions of the new edge-support equivalence
follow.  Finally, for every trunk index in that same range, replace the two
endpoint vertices and the exact edge by their proved prefix equalities in
$N$'s trunk-edge endpoint equation.  This proves the full trunk-endpoint
field.  All fields of normal globally maximum data have now been supplied,
and by construction the new height is $p$.
\end{proof}
